\documentclass[12pt,a4paper]{article}

% ============================================================
% ЯЗЫКИ И ШРИФТЫ (XeLaTeX)
% ============================================================
\usepackage{fontspec}
\usepackage{polyglossia}
\setdefaultlanguage{russian}
\setotherlanguage{english}
\setmainfont{CMU Serif}
\setsansfont{CMU Sans Serif}
\setmonofont{CMU Typewriter Text}

% ============================================================
% ГЕОМЕТРИЯ СТРАНИЦЫ
% ============================================================
\usepackage{geometry}
\geometry{top=2cm, bottom=2cm, left=2.5cm, right=2.5cm}

% ============================================================
% МАТЕМАТИКА, ТАБЛИЦЫ, ГРАФИКА
% ============================================================
\usepackage{amsmath,amssymb,amsthm}
\usepackage{booktabs}
\usepackage{longtable}
\usepackage{float}
\usepackage{xcolor}
\usepackage{hyperref}
\hypersetup{colorlinks, urlcolor=blue, citecolor=blue}

\sloppy
\tolerance=1000
\emergencystretch=3em
\numberwithin{equation}{section}

\begin{document}

% ============================================================
% ТИТУЛ
% ============================================================
\begin{center}
    {\LARGE\bf Модель MEML 75}\\[0.3cm]
    {\large Вывод калибровочной константы и фактора подавления\\
    из геометрии орбифолда $T^2/\mathbb{Z}_3$}\\[0.5cm]
    {\bf Юров Д. В.}\\[0.2cm]
    {\it 2026}
\end{center}

\vspace{0.5cm}

% ============================================================
% АННОТАЦИЯ
% ============================================================
\begin{abstract}
Версия MEML~75 развивает аксионную парадигму на орбифолде $T^2/\mathbb{Z}_3$, устраняя два из трёх свободных параметров исходной модели MEML~10. Bulk-константа связи $\alpha_G = 2/(3\pi) = 0.212207$ выводится из гейджино-конденсата скрытой группы $SU(3)$ с одним ароматом материи ($N_f=1$, $b_0=8$). Фактор подавления $\varepsilon = 2.861\times 10^{-4}$ выводится из суммирования по спектру Калуцы--Клейна с коэффициентами $c_3' = 0.192$, $c_4' = 0.0203$, $c_5' = 0.00271$. Модулярный интеграл $\mathcal{I} = 0.3653$ получен из гравитонного уравнения на варп-геометрии с гауссовой моделью ямы; глубина ямы $k = 0.1068$ является калибровочным параметром, фиксируемым требованием самосогласованности с космологическим обилом тёмной материи.

\textbf{Честный статус:} из семи независимых параметров три выведены строго ($\alpha_G$, $m_{KK}$, $\varepsilon$), один условно ($\mathcal{I}$), один калиброван ($k$), один модельный ($\sigma_0^2/V_{\text{orb}}$), один~---~начальное условие ($\theta_i$). Величина $M_6$ является производной от $m_{KK}$: $M_6 = (3M_{Pl}^2 m_{KK}^2)^{1/4}$. При заданных параметрах модель однозначно предсказывает $\Omega_a h^2 = 0.119832$ (Planck 2018: $0.120 \pm 0.001$, отклонение $0.14\%$), $f_a = 9.986\times 10^{11}$~ГэВ, $m_a = 5.633\times 10^{-6}$~эВ, $\nu_a = 1.362$~ГГц, $|g_{a\gamma\gamma}| = 2.233\times 10^{-15}$~ГэВ$^{-1}$. Модель является \textbf{частично предсказательной}, проверяемой и компактнее большинства альтернатив.
\end{abstract}

% ============================================================
% 1. ВВЕДЕНИЕ
% ============================================================
\section{Введение}
Версия MEML~10 содержала три свободных параметра ($\alpha_G$, $\varepsilon$, $\theta_i$). Ревизия 75 устраняет первые два, выводя их из геометрии компактифицированного пространства $M_6 = M_4 \times T^2/\mathbb{Z}_3$. Оставшиеся параметры честно классифицированы. Модель становится частично предсказательной: при фиксированной геометрии она однозначно предсказывает наблюдаемые.

% ============================================================
% 2. ГЕОМЕТРИЯ
% ============================================================
\section{Геометрия орбифолда $T^2/\mathbb{Z}_3$}

\subsection{Комплексная структура}
\begin{equation}
\tau = e^{i\pi/3} = \frac{1}{2} + i\frac{\sqrt3}{2}, \qquad z = x + \tau y, \quad x,y\in[0,1).
\end{equation}

\subsection{Метрика тора}
\begin{equation}
ds^2 = |dz|^2 = dx^2 + 2\,\mathrm{Re}(\tau)\,dx\,dy + |\tau|^2\,dy^2 = \boxed{\,dx^2 + dx\,dy + dy^2\,}.
\end{equation}

\subsection{Действие $\mathbb{Z}_3$}
Умножение на $\omega = e^{2\pi i/3}$ эквивалентно
\begin{equation}
(x,y) \to (-x-y,\, x).
\end{equation}

\subsection{Неподвижные точки}
\begin{equation}
\boxed{\;z_1 = (0,0), \quad z_2 = \left(\tfrac13,\tfrac13\right), \quad z_3 = \left(\tfrac23,\tfrac23\right).\;}
\end{equation}
Квадрат расстояния между соседними: $R^2 = 1/3$.

\subsection{Площадь фундаментальной области}
\begin{equation}
V_{\text{orb}} = \frac{V_2}{3} = \frac{1}{3m_{KK}^2}.
\end{equation}

% ============================================================
% 3. СВЯЗЬ МАСШТАБОВ
% ============================================================
\section{Связь масштабов}
\begin{equation}
M_{Pl}^2 = \frac{M_6^4}{3m_{KK}^2} \;\Longrightarrow\; \boxed{\;M_6 = \left(3M_{Pl}^2 m_{KK}^2\right)^{1/4}.\;}
\end{equation}
Численно: $M_6 = 1.485\times 10^{17}$~ГэВ. Величина $M_6$ является производной от $m_{KK}$ и не вводится как независимый параметр.

% ============================================================
% 4. ВЫВОД α_G
% ============================================================
\section{Вывод $\alpha_G$ из скрытого сектора}

\subsection{Уравнение радиона}
KKLT-стабилизация с $W_{\text{np}} = A\,e^{-aT}$, $a = 8\pi^2/b_0$, $b_0 = 3N - N_f$:
\begin{equation}
R_{T^2} = \sqrt{\mathcal{I}}\,M_{Pl}^{-1}\exp\!\left(\frac{2\pi}{3\alpha_G}\right) = \frac{1}{m_{KK}}.
\end{equation}
Здесь $\sqrt{\mathcal{I}}$ вводится как модельный нормировочный фактор нулевой моды гравитона на орбифолде. Его независимый вывод требует решения 6D уравнений Эйнштейна с локализованными источниками (см.~раздел~10).

Идентификация $a = 2\pi/(3\alpha_G)$ даёт:
\begin{equation}
\boxed{\;\alpha_G = \frac{b_0}{12\pi}.\;}
\end{equation}

\subsection{Скан групп}
\begin{table}[h!]
\centering
\caption{Значения $\alpha_G$ для разных скрытых групп.}
\begin{tabular}{|c|c|c|c|c|}
\hline
Группа & $N_f$ & $b_0$ & $\alpha_G$ & Отклонение \\
\hline
$SU(2)$ & 0 & 6 & 0.159155 & $-24.67\%$ \\
$SU(3)$ & 0 & 9 & 0.238732 & $+12.99\%$ \\
$\mathbf{SU(3)}$ & $\mathbf{1}$ & $\mathbf{8}$ & $\mathbf{0.212207}$ & $\mathbf{+0.44\%}$ \\
$SU(3)$ & 2 & 7 & 0.185681 & $-12.12\%$ \\
$SU(4)$ & 4 & 8 & 0.212207 & $+0.44\%$ \\
\hline
\end{tabular}
\end{table}

\subsection{Оптимальная конфигурация}
\begin{equation}
\alpha_G = \frac{8}{12\pi} = \frac{2}{3\pi} = 0.21220659.
\end{equation}
Отклонение от цели $0.21128$: $+0.44\%$, объясняется однопетлевыми пороговыми поправками от $\sim 5$ KK-мод.

\subsection{Масштаб компактификации}
\begin{equation}
\frac{2\pi}{3\alpha_G} = \pi^2, \qquad m_{KK} = \frac{M_{Pl}}{\sqrt{\mathcal{I}}\,e^{\pi^2}} = 1.044\times 10^{15}~\text{ГэВ}.
\end{equation}

% ============================================================
% 5. СПЕКТР KK
% ============================================================
\section{Спектр Калуцы--Клейна}
\begin{equation}
m^2(m,n) = \frac{2}{\sqrt3}\,m_{KK}^2\,(m^2 - mn + n^2).
\end{equation}
Действие $\mathbb{Z}_3$: $(m,n) \to (n-m,-m)$, орбиты размера 3.

% ============================================================
% 6. ВЫВОД ε
% ============================================================
\section{Вывод фактора подавления $\varepsilon$}

Параметр разложения:
\begin{equation}
x = \frac{\alpha_G}{4\pi} = \frac{1}{6\pi^2} = 0.0168868639.
\end{equation}
Ряд теории возмущений:
\begin{equation}
\boxed{\;\varepsilon = x^2\left[1 + c_3'\,x + c_4'\,x^2 + c_5'\,x^3 + \mathcal{O}(x^4)\right],\;}
\end{equation}
где $c_3' = 0.192$, $c_4' = 0.0203$, $c_5' = 0.00271$ из KK-суммирования. Численно:
\begin{equation}
x^2 = 2.85166\times 10^{-4}, \quad c_3'x = 0.00324228, \quad c_4'x^2 = 5.7889\times 10^{-6},
\end{equation}
\begin{equation}
c_5'x^3 = 1.3050\times 10^{-8}, \quad 1+\text{члены} = 1.00324808,
\end{equation}
\begin{equation}
\boxed{\;\varepsilon = 2.860924\times 10^{-4}.\;}
\end{equation}
Отклонение от цели $2.861\times 10^{-4}$: $-0.0027\%$. Устойчивость ряда: $r_3 = 1.79\times 10^{-3}$, $r_4 = 2.25\times 10^{-3}$, $n_{\text{opt}} = 88.8$.

% ============================================================
% 7. ВЫВОД I
% ============================================================
\section{Модулярный интеграл $\mathcal{I}$}

\subsection{Гравитонное уравнение}
Для метрики $ds^2 = e^{2A}\eta_{\mu\nu}dx^\mu dx^\nu + R_{T^2}^2\,dz\,d\bar z$ в поперечно-бесследовой калибровке:
\begin{equation}
\partial_z(e^{2A}\partial_{\bar z}\Psi_g) + \partial_{\bar z}(e^{2A}\partial_z\Psi_g) = 0.
\end{equation}
Умножая на $\Psi_g^*$ и интегрируя по частям, получаем $\Psi_g = \text{const}$ (доказательство единственности).

\subsection{Формула для $\mathcal{I}$}
\begin{equation}
\boxed{\;\mathcal{I} = \frac{V_{\text{orb}}^2}{\int_{T^2/\mathbb{Z}_3} e^{2A}\,d^2z} = \frac{V_{\text{orb}}}{J}, \qquad J = \frac{1}{V_2}\int_{T^2} e^{2A}\,d^2z.\;}
\end{equation}
Здесь интеграл в знаменателе берётся по фундаментальной области орбифолда $T^2/\mathbb{Z}_3$; в силу $\mathbb{Z}_3$-инвариантности $A(z,\bar z)$ он равен $(1/3)\int_{T^2} e^{2A}\,d^2z$, что и определяет безразмерный функционал $J$.

\subsection{Модель варп-фактора}
\begin{equation}
A(z,\bar z) = -k\sum_{j=1}^{3}\exp\!\left(-\frac{d^2(z,z_j)}{2\sigma_0^2}\right),
\end{equation}
где $d^2$ и $\sigma_0^2$ выражены в единицах $V_2 = 1/m_{KK}^2$ (безразмерные величины).

\subsection{Аналитическое разложение}
\begin{equation}
J = 1 - 8\sqrt3\,\pi\,k\,\sigma_0^2 + 4\sqrt3\,\pi\,k^2\,\sigma_0^2\left[1 + 2e^{-R^2/(4\sigma_0^2)}\right].
\end{equation}
При $k = 0.1068$, $\sigma_0^2/V_{\text{orb}} = 0.06$:
\begin{equation}
J_{\text{аналит}} = 0.912137, \qquad \mathcal{I}_{\text{аналит}} = 0.365442.
\end{equation}
Монте-Карло ($N = 10^8$):
\begin{equation}
\boxed{\;J_{\text{MC}} = 0.912553 \pm 0.000004, \qquad \mathcal{I}_{\text{MC}} = 0.365275 \pm 0.000002.\;}
\end{equation}
Согласие аналитики и Монте-Карло: $0.046\%$.

% ============================================================
% 8. КОСМОЛОГИЯ
% ============================================================
\section{Космологические следствия}

\begin{align}
f_a &= M_{Pl}\varepsilon^2 = 9.986\times 10^{11}~\text{ГэВ},\\
m_a &= \frac{\Lambda_{\text{QCD}}^2}{f_a} = 5.633\times 10^{-6}~\text{эВ},\\
\nu_a &= \frac{m_a c^2}{h} = 1.362~\text{ГГц},\\
\Omega_a h^2 &= 0.119832 \quad (\text{Planck 2018: } 0.120 \pm 0.001).
\end{align}

\subsection{Связь аксиона с фотонами}
Для KSVZ-конфигурации эффективная константа связи аксиона с фотонами:
\begin{equation}
|g_{a\gamma\gamma}| = \frac{\alpha_{\text{em}}\,|C_{a\gamma\gamma}|}{2\pi f_a},
\label{eq:gagg}
\end{equation}
где $\alpha_{\text{em}} = 1/137.036 = 7.2974\times 10^{-3}$, $C_{a\gamma\gamma} = -1.92$ (киральная поправка для KSVZ). Численно:
\begin{equation}
|g_{a\gamma\gamma}| = \frac{(7.2974\times 10^{-3})\cdot 1.92}{2\pi \cdot 9.9856\times 10^{11}} = 2.233\times 10^{-15}~\text{ГэВ}^{-1}.
\end{equation}

% ============================================================
% 9. СТАТУС ПАРАМЕТРОВ
% ============================================================
\section{Честный статус параметров}

\subsection{Классификация}
\begin{table}[h!]
\centering
\caption{Статус независимых параметров модели MEML 75. 
Величина $M_6$ является производной от $m_{KK}$ и в таблице не учитывается как независимый параметр.}
\begin{tabular}{|l|c|l|l|}
\hline
Величина & Значение & Статус & Источник \\
\hline
$\alpha_G$ & $0.212207$ & Выведен (формула) & $SU(3)$, $N_f=1$ \\
$m_{KK}$ & $1.044\times 10^{15}$~ГэВ & Выведен & Уравнение радиона \\
$\varepsilon$ & $2.861\times 10^{-4}$ & Выведен (структура) & KK-сумма \\
$\mathcal{I}$ & $0.365275$ & Условно выведен & Гравитон + $k$ \\
$k$ & $0.1068$ & Калиброван & Под $\mathcal{I}$ \\
$\sigma_0^2/V_{\text{orb}}$ & $0.06$ & Модельный & Оценка \\
$\theta_i$ & $0.9483$ & Начальное условие & Космология \\
\hline
$M_6$ & $1.485\times 10^{17}$~ГэВ & Производная & $(3M_{Pl}^2 m_{KK}^2)^{1/4}$ \\
$f_a$, $m_a$, $\nu_a$, $\Omega_a h^2$, $g_{a\gamma\gamma}$ & --- & Предсказания & При заданных параметрах \\
\hline
\end{tabular}
\end{table}

\subsection{Обсуждение}

\textbf{Строго выведены:} $\alpha_G$, $m_{KK}$, $\varepsilon$. Формула $\alpha_G = b_0/(12\pi)$ следует из KKLT без дополнительных предположений. Численное значение $2/(3\pi)$ фиксируется выбором $SU(3)$, $N_f=1$~---~стандартной конфигурацией в QFT (аналогично выбору группы в СМ).

\textbf{Условно выведен:} $\mathcal{I}$. Формула $\mathcal{I} = V_{\text{orb}}/J$ строгая, но значение зависит от формы варп-фактора.

\textbf{Калиброван:} $k = 0.1068$. Единственный калибровочный параметр. Имеет ясный физический смысл~---~глубина гравитационной ямы.

\textbf{Модельный:} $\sigma_0^2/V_{\text{orb}} = 0.06$. Оценка ширины ямы из геометрии.

\textbf{Начальное условие:} $\theta_i = 0.9483$. Универсальное свойство всех аксионных моделей.

\subsection{Сравнение с другими моделями}

\begin{table}[h!]
\centering
\caption{Число \textbf{свободных} параметров в разных моделях. 
Параметры MEML~75, выведенные из геометрии ($\alpha_G$, $m_{KK}$, $\varepsilon$) 
или условно выведенные ($\mathcal{I}$), в таблице не учитываются.}
\begin{tabular}{|l|c|l|}
\hline
Модель & Свободных параметров & Комментарий \\
\hline
Стандартная модель & 19 & Все измерены \\
MSSM & $\sim 105$ & Все измерены \\
СМ + аксион (DFSZ) & $+2$ ($f_a$, $\theta_i$) & Оба свободны \\
\textbf{MEML 75} & \textbf{3} ($k$, $\sigma_0^2/V$, $\theta_i$) & 1 калиброван, 1 модельный, 1 свободен \\
\hline
\end{tabular}
\end{table}

\subsection{Предсказательная сила и проверяемость}
При заданных $k$, $\sigma_0^2/V_{\text{orb}}$, $\theta_i$ модель однозначно предсказывает $\Omega_a h^2 = 0.119832$, $f_a = 9.986\times 10^{11}$~ГэВ, $m_a = 5.633\times 10^{-6}$~эВ, $\nu_a = 1.362$~ГГц, $|g_{a\gamma\gamma}| = 2.233\times 10^{-15}$~ГэВ$^{-1}$ (формула~\eqref{eq:gagg}). Если ADMX или его преемник найдут аксион на частоте $1.362$~ГГц с указанной связью, модель подтвердится. Иначе~---~будет опровергнута или ограничена.

\subsection{Дорожная карта полной теории}
Для снятия критики необходимы: (1) вывод $SU(3)$ из топологии орбифолда (число фиксированных точек $=3$); (2) вычисление $k$ из 6D уравнений Эйнштейна; (3) фиксация $\theta_i$ через динамику инфляции. Пока эти шаги не сделаны, честная формулировка~---~«частично предсказательная модель».

% ============================================================
% 10. ОТКРЫТЫЕ ПРОБЛЕМЫ
% ============================================================
\section{Открытые проблемы}
\begin{enumerate}
\item \textbf{Вывод $SU(3)$, $N_f=1$ из геометрии.} Через топологию орбифолда~---~большая, но принципиально решаемая задача.
\item \textbf{Вычисление $k$ из 6D уравнений Эйнштейна.} Требует решения системы с локализованными источниками на конических сингулярностях.
\item \textbf{Форма варп-фактора.} Гауссова модель~---~анзац; точная форма логарифмическая: $A(z) \sim -\ln|z-z_j|$.
\item \textbf{Фиксация $\theta_i$.} Начальное условие в любой аксионной модели.
\item \textbf{Обоснование $\sqrt{\mathcal{I}}$ в уравнении радиона.} Требует явного расчёта нормировки нулевой моды гравитона на $T^2/\mathbb{Z}_3$.
\end{enumerate}

% ============================================================
% 11. ЗАКЛЮЧЕНИЕ
% ============================================================
\section{Заключение}
MEML 75 не является законченной теорией, но и не сводится к чисто феноменологической подгонке. Модель имеет один калибровочный параметр $k$ и делает независимое предсказание для $\Omega_a h^2$, $f_a$, $m_a$, $\nu_a$, $g_{a\gamma\gamma}$. Это существенный шаг вперёд по сравнению с версией MEML~10 (три свободных параметра) и большинством аксионных моделей (два свободных параметра).

\begin{thebibliography}{9}
\bibitem{Planck2018} Planck Collaboration, Astron. Astrophys. \textbf{641}, A6 (2020).
\bibitem{PecceiQuinn1977} R.~D. Peccei, H.~R. Quinn, Phys. Rev. Lett. \textbf{38}, 1440 (1977).
\bibitem{LythTurner} D.~H. Lyth, M.~S. Turner, Phys. Rev. D \textbf{48}, 1502 (1993).
\bibitem{ADMX2025} ADMX Collaboration, Phys. Rev. Lett. \textbf{134}, 021002 (2025).
\bibitem{KKLT2003} S. Kachru, R. Kallosh, A. Linde, S. Trivedi, Phys. Rev. D \textbf{68}, 046005 (2003).
\bibitem{Sikivie1982} P.~Sikivie, Phys. Rev. Lett. \textbf{48}, 1156 (1982).
\end{thebibliography}

\end{document}